\begin{abstract}

Model-heterogeneous federated learning (HFL) allows clients with different architectures to collaborate, but corruption-robust HFL usually treats corruption as nuisance variation. We study \textbf{class--corruption entanglement (CLE)}, a controlled setting in which client-specific associations between task classes and corruption operators can make corruption a directional class cue. We evaluate the same semantic source under multiple operators and introduce \textbf{Directional Shortcut Alignment (DSA)}, which measures whether probability mass moves toward the classes associated with the applied operator during training. Under explicit assumptions about semantic preservation, pre-registered bindings, and evaluation isolation, DSA measures an operator-response contrast along the training binding; a proxy non-identifiability result shows why pseudo-environment improvements alone cannot certify reduced CLE dependence. A seed-0 matched HFL-versus-Local factorial finds a predominantly local-first formation pattern, motivating a local intervention. We combine a taxonomy-assisted coarse \textbf{Public Environment Witness (PEW)} with \textbf{Balanced Environment Risk (BER)} to compress class-conditional pseudo-environment support imbalance without changing communication. On a held-out binding map over three matched training seeds, PEW+BER attains lower pooled DSA and higher pooled task utility than ERM, CVaR-DRO, and PEW+GroupDRO. Seed-0 experiments also reproduce DSA suppression on another binding map and under a second protocol-matched HFL communication base, where the utility gate fails. The method is taxonomy-assisted, and shortcut suppression does not imply uniform utility gains.

\end{abstract}
